\section*{ВВЕДЕНИЕ}
\addcontentsline{toc}{section}{ВВЕДЕНИЕ}

Компьютерная графика представляет собой совокупность методов и средств преобразования информации в графическое представление с использованием вычислительной техники. В области театрального искусства цифровые технологии позволяют значительно упростить и ускорить процесс проектирования оформления сцены, визуализации декораций, настройки освещения и оценки восприятия постановки зрителем ещё до начала физического изготовления декораций и монтажа светового оборудования.

Традиционное проектирование театральных декораций опирается на эскизы, макеты и чертежи, что требует значительных временных и материальных затрат. Использование трёхмерного моделирования даёт возможность создавать виртуальные макеты сцены, учитывающие геометрические параметры площадки, расположение порталов, глубину сцены, технические ограничения конкретного театра, а также взаимодействие света с поверхностями декораций. Такие инструменты широко применяются в современной сценографии и позволяют проводить предвизуализацию (previsualization) спектакля.

Актуальность темы обусловлена растущей потребностью театров в цифровых средствах проектирования. Системы автоматизированного проектирования и специализированные программные комплексы (такие как Vectorworks Spotlight, WYSIWYG, Capture, а также универсальные пакеты Blender, Unreal Engine) активно используются для моделирования декораций, светового оформления и координации работы постановочных цехов. Вместе с тем для учебных и исследовательских целей, а также для небольших постановок востребованы относительно простые и специализированные программные комплексы, ориентированные именно на моделирование оформления театральной сцены с интерактивным управлением камерой, источниками света и параметрами объектов.

Целью курсового проекта является разработка программного комплекса моделирования оформления театральной сцены, позволяющего создавать и редактировать трёхмерные композиции из декорационных элементов, настраивать параметры освещения и визуализировать сцену с учётом удаления невидимых поверхностей, закраски и построения теней.

Для достижения поставленной цели необходимо решить следующие задачи:

\begin{itemize}
	\item провести анализ предметной области и существующих методов и алгоритмов компьютерной графики, применяемых при моделировании трёхмерных сцен и театрального оформления;
	\item выбрать наиболее подходящие модели представления объектов, алгоритмы удаления невидимых линий и поверхностей, закраски и освещения;
	\item спроектировать архитектуру программного комплекса и его графический интерфейс пользователя;
	\item выбрать средства программной реализации;
	\item разработать программное обеспечение, реализующее выбранные алгоритмы и структуры данных;
	\item провести экспериментальное исследование временных характеристик разработанного программного комплекса.
\end{itemize}

Решение перечисленных задач описывается в соответствующих разделах расчётно-пояснительной записки.

\newpage